\documentclass[11pt]{article}
\usepackage[margin=0.9in]{geometry}
\usepackage{fontspec}
\usepackage{amsmath,amssymb,amsthm}
\usepackage{xurl}
\begin{document}\sloppy
\begin{center}
{\Large GROK H0033: first-step cube 2 of \texttt{a40ff54e580f5d0a} at $L=9$ closes UNSAT through all 210 subcubes, and 158 of cube 0's 210 subcubes are certified}\\[4pt]
2026-10-06 00:33 America/Los\_Angeles \quad branch \texttt{main}, tip \texttt{c5265be4} \quad automation 41774ea2\\[2pt]
\textbf{P versus NP: OPEN.}
\end{center}

\section*{Peer delta since H0012 (\texttt{b522548d})}
Checked 00:34\,PT. Two commits landed. \texttt{86f9a9d0} (00:20\,PT) adds \texttt{research/explorations/monotone-slice/2026-10-06\_0040\_TWOSATDEC14.md} and \texttt{results/raw/monotone\_slice/TWOSATDEC14\_q19.txt}. \texttt{c5265be4} (00:32\,PT) adds \texttt{research/agents/RED/2026-10-06\_RED-0007\_trace\_noninvariance.md}. Neither touches \texttt{research/main/} or \texttt{research/agents/grok/}. Their contents are unread here, and nothing below depends on them. Root \texttt{main.tex} is not touched by this note.

\section*{1. Literature}
None needed. The next claim is again about our own finite search: whether the 210 split pieces of one more first-step cube are all refuted.

\section*{2. Claims recorded}
\textbf{Claim CUBE2 (COMPUTATIONAL, solver-trusted).} First-step cube 2, the step $(w,s_a,s_b,g)=(0,0,1,6)$, of the tight $L=9$, $r=2$, $n=6$ search for \texttt{a40ff54e580f5d0a} is UNSAT. Phase B of \texttt{split\_queue.sh} wrote \texttt{CUBE-DONE 2 00:34:04}. Its 210 subcube lines name 210 distinct indices and 210 distinct step-2 tuples, all UNSAT under CaDiCaL, with no SAT and no TIMEOUT line anywhere in the split log. Solve times: median 11.6\,s, slowest 376.5\,s, sum 6919.4\,s, 135 over 1\,s. These subcubes carry no DRAT certificate.

\textbf{Claim DRAT158 (COMPUTATIONAL).} 158 of cube 0's 210 subcubes carry drat-trim-verified refutations. The 7 new ones are 190, 103, 53, 199, 127, 158 and 131. That leaves 52, and subcube 191 was being checked at 00:34\,PT.

\textbf{Correction (arithmetic, not a claim change).} H0012 gave cube 1's median solve time as 11.6\,s. Recomputed over its 210 lines, the two middle values are 11.2\,s and 11.6\,s, so the median is 11.4\,s; 11.6\,s was the upper median. Its slowest time (305.6\,s), sum (5499.6\,s) and count over 1\,s (135) recompute exactly. CUBE1 itself is unaffected.

\section*{3. Smallest honest test}
\begin{itemize}
\item \textbf{Split-log audit} (\texttt{/workspace/scratch\_h0033\_1006/experiments/cube\_check.py}, 00:35\,PT). Parses every subcube line of \texttt{a40\_L9\_split.out}. Cubes 0, 1 and 2 each have 210 lines, 210 distinct indices, 210 distinct step-2 tuples and 210 UNSAT. No line outside the pattern mentions SAT or TIMEOUT. For every subcube $j$, cubes 1, 2 and the running cube 4 fix the same step-2 tuple as cube 0's subcube $j$, so all of them use the 210-way map whose cover argument CUBE0 already relies on.
\item \textbf{Index control} (inline Python, 00:35\,PT). \texttt{regsat\_fix.cubes(6)} and \texttt{regsat\_cube.cubes(6)} are equal, 150 entries each. Entry 2 is $(0,0,1,6)$ and entry 4 is $(0,0,1,14)$, matching the tuples printed on the phase-B lines. Cube 3 was UNSAT in the first-step runner, so phase B correctly skips from 2 to 4.
\item \textbf{Recount} (\texttt{/workspace/scratch\_h2333\_1005/experiments/recount.py}, rerun 00:35\,PT). 75 trivial-UP lines plus 6 more give 81. The 129-entry queue is disjoint from them and the two partition $\{0,\dots,209\}$. The H2301 batch has written 77 lines in queue order, all \texttt{s VERIFIED}, with 77 distinct CNF hashes, giving 158. Solve times 7.8 to 64.6\,s, check times 7.9 to 127.3\,s, check/solve 0.90 to 2.58.
\item \textbf{Coverage control} (runner log, 00:35\,PT). Cubes 0 to 74 each have exactly one result line: 42 UNSAT, 33 TIMEOUT, 0 SAT. With cubes 0, 1 and 2 closed by splitting, 45 of the 75 kept $w=0$ cubes are closed and 30 remain open. Phase B is on cube 4 and had 47 of its subcubes UNSAT, none SAT, at 00:35\,PT.
\end{itemize}

\section*{4. What this changes}
One fewer open first-step cube, from 31 to 30. The bounds stay $9\le L(\texttt{a40ff54e580f5d0a},2)\le 11$ (COMPUTATIONAL). The lower bound would rise to 10 only if all 30 open cubes also close. The largest confirmed price stays one.

\section*{Effect on the standing candidate}
No claim is falsified and no charged resource improves. The time-state-representation surface stays CANDIDATE. The SAT-to-charged-model compiler is OPEN, and so is the product-span obstruction. These results cover one fixed function and say nothing about lower bounds for any family.

\section*{Next smallest step}
Let phase B continue through cube 4 and the remaining timed-out cubes, and let the DRAT batch finish cube 0's last 52 subcubes. Once \texttt{split\_queue.sh} exits, check that the guard's \texttt{CATCHUP} lines name 72, 74 and, only if phase B skipped it, 69.

\bigskip
\noindent\textbf{Status sentence.} P versus NP remains OPEN: the SAT-to-charged-model compiler and the product-span obstruction are open load-bearing arrows. This pass records that first-step cube 2 of \texttt{a40ff54e580f5d0a} at tight $L=9$ is UNSAT through all 210 of its subcubes (COMPUTATIONAL, solver-trusted, cube index checked against the runner's list) and that 158 of cube 0's 210 subcubes now carry drat-trim-verified refutations (COMPUTATIONAL). Among the 75 kept first-step cubes, 45 are closed and 30 are open, with 0 SAT, and $9\le L(f,2)\le 11$ is unchanged. No claim is falsified, the largest confirmed price stays one, no charged resource improves, and this is not a proof of anything about P versus NP.
\end{document}
